% Part: first-order-logic
% Chapter: natural-deduction
% Section: proving-things-quant

\documentclass[../../../include/open-logic-section]{subfiles}

\begin{document}

\olfileid{fol}{ntd}{prq}

\olsection{\usetoken{P}{derivation} met kwantoren}

\begin{ex}
Bij kwantoren moeten we erop letten dat we de
eigenvariabelevoorwaarde niet schenden. Deze voorwaarde bepaalt dat een
nieuw gekozen symbool niet al op de verboden plaatsen mag voorkomen.
Soms moeten we daarom de volgorde van bepaalde inferenties aanpassen.
Het helpt om de regels met deze voorwaarde het eerst te behandelen. In
het voltooide bewijs staan die regels lager.

We bekijken hoe we !!a{derivation} maken van de !!{formula}
$\lexists[x][\lnot !A(x)] \lif \lnot \lforall[x][!A(x)]$.
Zoals altijd beginnen we met:
\begin{prooftree}
\AxiomC{}
\UnaryInfC{$\lexists[x][\lnot !A(x)]\lif \lnot \lforall[x][!A(x)]$}
\end{prooftree}
Eerst schrijven we op wat we nodig hebben om de laatste stap met de
regel \Intro{\lif} te verantwoorden.
\begin{prooftree}
\AxiomC{$\Discharge{\lexists[x][\lnot !A(x)]}{1}$}
\DeduceC{$\lnot \lforall[x][!A(x)]$}
\DischargeRule{\Intro{\lif}}{1}
\UnaryInfC{$\lexists[x][\lnot !A(x)]\lif \lnot \lforall[x][!A(x)]$}
\end{prooftree}
Er is geen duidelijke regel die we meteen op
$\lnot \lforall[x][!A(x)]$ kunnen toepassen. Daarom zetten we de
!!{derivation} zo op dat we \Elim{\lexists} kunnen gebruiken. Nu
moeten we op de eigenvariabelevoorwaarde letten. We kiezen een
constante die niet voorkomt in $\lexists[x][!A(x)]$ of in een aanname
waarvan deze formule afhangt. (Er komen hier geen !!{constant}s voor,
dus elke keuze is goed.)
\begin{prooftree}
\AxiomC{$\Discharge{\lexists[x][\lnot !A(x)]}{1}$}
\AxiomC{$\Discharge{\lnot !A(a)}{2}$}
\DeduceC{$\lnot \lforall[x][!A(x)]$}
\DischargeRule{\Elim{\lexists}}{2}
\BinaryInfC{$\lnot \lforall[x][!A(x)]$}
\DischargeRule{\Intro{\lif}}{1}
\UnaryInfC{$\lexists[x][\lnot !A(x)] \lif \lnot \lforall[x][!A(x)]$}
\end{prooftree}
Om $\lnot \lforall[x][!A(x)]$ af te leiden, proberen we
\Intro{\lnot}. Daarvoor moeten we een tegenspraak afleiden. We mogen
$\lforall[x][!A(x)]$ daarbij als extra aanname gebruiken. De
tegenspraak mag ook afhangen van de aanname $\lnot !A(a)$. De
\Elim{\lexists}-inferentie zal die aanname later opheffen. We zetten
de afleiding zo op:
\begin{prooftree}
\AxiomC{$\Discharge{\lexists[x][\lnot !A(x)]}{1}$}
\AxiomC{$\Discharge{\lnot !A(a)}{2}, \Discharge{\lforall[x][!A(x)]}{3}$}
\DeduceC{$\lfalse$}
\DischargeRule{\Intro{\lnot}}{3}
\UnaryInfC{$\lnot \lforall[x][!A(x)]$}
\DischargeRule{\Elim{\lexists}}{2}
\BinaryInfC{$\lnot \lforall[x][!A(x)]$}
\DischargeRule{\Intro{\lif}}{1}
\UnaryInfC{$\lexists[x][\lnot !A(x)]\lif \lnot \lforall[x][!A(x)]$}
\end{prooftree}
We zijn bijna bij een tegenspraak. De gemakkelijkste regel is
\Elim{\lforall}; daarvoor geldt geen eigenvariabelevoorwaarde. We
mogen de universeel gekwantificeerde~$x$ door elke term vervangen.
Daarom gebruiken we opnieuw~$a$, zodat we een tegenspraak krijgen.
\begin{prooftree}
\AxiomC{$\Discharge{\lexists[x][\lnot !A(x)]}{1}$}
\AxiomC{$\Discharge{\lnot !A(a)}{2}$}
\AxiomC{$\Discharge{\lforall[x][!A(x)]}{3}$}
\RightLabel{\Elim{\lforall}}
\UnaryInfC{$!A(a)$}
\RightLabel{\Elim{\lnot}}
\BinaryInfC{$\lfalse$}
\DischargeRule{\Intro{\lnot}}{3}
\UnaryInfC{$\lnot \lforall[x][!A(x)]$}
\DischargeRule{\Elim{\lexists}}{2}
\BinaryInfC{$\lnot \lforall[x][!A(x)]$}
\DischargeRule{\Intro{\lif}}{1}
\UnaryInfC{$\lexists[x][\lnot !A(x)]\lif \lnot \lforall[x][!A(x)]$}
\end{prooftree}

Vooral bij kwantoren moeten we nu nog eens controleren of de
eigenvariabelevoorwaarde nergens is geschonden. De enige toegepaste
regel met die voorwaarde was \Elim{\exists}. De eigenvariabele~$a$
komt niet voor in een aanname waarvan de betreffende formule afhangt.
Dit is dus een correcte !!{derivation}.
\end{ex}

\begin{ex}
Soms leiden we !!a{formula} af uit andere !!{formula}s. Dan kunnen er
aannamen overblijven die niet zijn opgeheven. We moeten zowel die
aannamen als het einddoel blijven bijhouden.

We bekijken hoe we !!a{derivation} maken van de !!{formula}
$\lexists[x][!C(x,b)]$ uit de aannamen $\lexists[x][(!A(x) 
\land !B(x))]$ en $\lforall[x][(!B(x) \lif !C(x,b))]$.
Zoals altijd schrijven we de conclusie onderaan.
\begin{prooftree}
\AxiomC{}
\UnaryInfC{$\lexists[x][!C(x,b)]$}
\end{prooftree}

We hebben twee uitgangsformules. Om de eerste te gebruiken, zoeken we
!!a{derivation} van $\lexists[x][!C(x, b)]$ uit
$\lexists[x][(!A(x) \land !B(x))]$. Daarvoor gebruiken we
\Elim{\lexists}. Bij die regel hoort een eigenvariabelevoorwaarde, dus
we passen hem als eerste toe. Dat geeft:
\begin{prooftree}
\AxiomC{$\lexists[x][(!A(x) \land !B(x))]$}
\AxiomC{$\Discharge{!A(a) \land !B(a)}{1}$}
\DeduceC{$\lexists[x][!C(x,b)]$}
\DischargeRule{\Elim{\lexists}}{1}
\BinaryInfC{$\lexists[x][!C(x,b)]$}
\end{prooftree}
In beide aannamen komt~$!B$ voor. Daarom is het nuttig nu
\Elim{\land} toe te passen en~$!B(a)$ apart te verkrijgen.
\begin{prooftree}
\AxiomC{$\lexists[x][(!A(x) \land !B(x)])$}
\AxiomC{$\Discharge{!A(a) 
\land !B(a)}{1}$}
\RightLabel{\Elim{\land}}
\UnaryInfC{$!B(a)$}
\DeduceC{$\lexists[x][!C(x,b)]$}
\DischargeRule{\Elim{\lexists}}{1}
\BinaryInfC{$\lexists[x][!C(x,b)]$}
\end{prooftree}

De tweede aanname is~$\lforall[x][(!B(x) \lif !C(x,b))]$. Hier geldt
geen eigenvariabelevoorwaarde. Met \Elim{\lforall} mogen we dus voor
$x$ de !!{constant}~$a$ invullen. Zo krijgen we
$!B(a) \lif !C(a, b)$. We hebben nu zowel $!B(a) \lif !C(a,b)$ als
$!B(a)$. De volgende stap is eenvoudig: pas \Elim{\lif} toe.
\begin{prooftree}
\AxiomC{$\lexists[x][(!A(x) \land !B(x))]$}
\AxiomC{$\lforall[x][(!B(x) \lif !C(x,b))]$}
\RightLabel{\Elim{\lforall}}
\UnaryInfC{$!B(a) \lif !C(a,b)$}
\AxiomC{$\Discharge{!A(a) 
\land !B(a)}{1}$}
\RightLabel{\Elim{\land}}
\UnaryInfC{$!B(a)$}
\RightLabel{\Elim{\lif}}
\BinaryInfC{$!C(a,b)$}
\DeduceC{$\lexists[x][!C(x,b)]$}
\DischargeRule{\Elim{\lexists}}{1}
\insertBetweenHyps{\hspace{-5em}}
\BinaryInfC{$\lexists[x][!C(x,b)]$}
\end{prooftree}
We zijn er bijna!{} Nog één toepassing van \Intro{\lexists} en we
hebben ons doel bereikt.
\begin{prooftree}
\AxiomC{$\lexists[x][(!A(x) \land !B(x))]$}
\AxiomC{$\lforall[x][(!B(x) \lif !C(x,b))]$}
\RightLabel{\Elim{\lforall}}
\UnaryInfC{$!B(a) \lif !C(a,b)$}
\AxiomC{$\Discharge{!A(a) 
\land !B(a)}{1}$}
\RightLabel{\Elim{\land}}
\UnaryInfC{$!B(a)$}
\RightLabel{\Elim{\lif}}
\BinaryInfC{$!C(a,b)$}
\RightLabel{\Intro{\lexists}}
\UnaryInfC{$\lexists[x][!C(x,b)]$}
\DischargeRule{\Elim{\lexists}}{1}
\insertBetweenHyps{\hspace{-5em}}
\BinaryInfC{$\lexists[x][!C(x,b)]$}
\end{prooftree}
Bij elke stap hebben we gecontroleerd dat de
eigenvariabelevoorwaarden niet werden geschonden. Daarom weten we dat
dit een correcte !!{derivation} is.
\end{ex}

\begin{ex}
Geef !!a{derivation} van de !!{formula}
$\lnot\lforall[x][!A(x)]$ uit de aannamen $\lforall[x][!A(x)] 
\lif \lexists[y][!B(y)]$ en $\lnot\lexists[y][!B(y)]$.
Zoals altijd schrijven we de beoogde !!{formula} onderaan.
\begin{prooftree}
\AxiomC{}
\UnaryInfC{$\lnot\lforall[x][!A(x)]$}
\end{prooftree}
De laatste regel van de !!{derivation} is een ontkenning. Daarom
proberen we \Intro{\lnot}. We moeten dan uitzoeken hoe we een
tegenspraak kunnen !!{derive}.
\begin{prooftree}
\AxiomC{$\Discharge{\lforall[x][!A(x)]}{1}$}
\DeduceC{$\lfalse$}
\DischargeRule{\Intro{\lnot}}{1}
\UnaryInfC{$\lnot\lforall[x][!A(x)]$}
\end{prooftree}
Tot zover gaat het goed. We kunnen \Elim{\lforall} gebruiken, maar het
is niet duidelijk of dat ons dichter bij het doel brengt. Daarom
gebruiken we een van de aannamen.
$\lforall[x][!A(x)] \lif \lexists[y][!B(y)]$ maakt samen met
$\lforall[x][!A(x)]$ een toepassing van \Elim{\lif} mogelijk.
\begin{prooftree}
\AxiomC{$\lforall[x][!A(x)] \lif \lexists[y][!B(y)]$}
\AxiomC{$\Discharge{\lforall[x][!A(x)]}{1}$}
\RightLabel{\Elim{\lif}}
\BinaryInfC{$\lexists[y][!B(y)]$}
\DeduceC{$\lfalse$}
\DischargeRule{\Intro{\lnot}}{1}
\UnaryInfC{$\lnot\lforall[x][!A(x)]$}
\end{prooftree}
Er blijft nog één aanname over. Daarmee bereiken we via
\Elim{\lnot} een tegenspraak.
\begin{prooftree}
\AxiomC{$\lnot\lexists[y][!B(y)]$}
\AxiomC{$\lforall[x][!A(x)] \lif \lexists[y][!B(y)]$}
\AxiomC{$\Discharge{\lforall[x][!A(x)]}{1}$}
\RightLabel{\Elim{\lif}}
\BinaryInfC{$\lexists[y][!B(y)]$}
\RightLabel{\Elim{\lnot}}
\BinaryInfC{$\lfalse$}
\DischargeRule{\Intro{\lnot}}{1}
\UnaryInfC{$\lnot\lforall[x][!A(x)]$}
\end{prooftree}
\end{ex}

\begin{prob}
Geef !!{derivation}s die het volgende laten zien:
\begin{enumerate}
\item $\Proves (\lforall[x][!A(x)] \land \lforall[y][!B(y)]) \lif
\lforall[z][(!A(z) \land !B(z))]$.
\item $\Proves (\lexists[x][!A(x)] \lor \lexists[y][!B(y)]) \lif
\lexists[z][(!A(z) \lor !B(z))]$.
\item $\lforall[x][(!A(x) \lif !B)] \Proves \lexists[y][!A(y)] \lif !B$.
\item $\lforall[x][\lnot !A(x)] \Proves \lnot\lexists[x][!A(x)]$.
\item $\Proves \lnot\lexists[x][!A(x)] \lif \lforall[x][\lnot !A(x)]$.
\item $\Proves \lnot\lexists[x][\lforall[y][((!A(x,y) \lif \lnot
!A(y,y)) \land (\lnot !A(y,y) \lif !A(x,y)))]]$.
\end{enumerate}
\end{prob}

\begin{prob}
Geef !!{derivation}s die het volgende laten zien:
\begin{enumerate}
\item $\Proves \lnot\lforall[x][!A(x)] \lif \lexists[x][\lnot!A(x)]$.
\item $(\lforall[x][!A(x)] \lif !B) \Proves \lexists[y][(!A(y) \lif !B)]$.
\item $\Proves \lexists[x][(!A(x) \lif \lforall[y][!A(y)])]$.
\end{enumerate}
(Hiervoor is steeds de regel $\FalseCl$ nodig.)
\end{prob}
\end{document}
